\documentclass[10pt,a4paper]{amsart}
\usepackage[margin=19mm]{geometry}
\usepackage{amsmath,amssymb,mathtools}
\usepackage[scr=rsfs,cal=euler]{mathalfa}
\usepackage{xfrac}
\usepackage[unicode,hidelinks,bookmarks=false]{hyperref}
\newcommand{\ie}{i.e.\ }
\newcommand{\cf}{cf.\ }
\newcommand{\resp}{respectively\ }
\newcommand{\Subsection}{\subsection}
\providecommand{\nobreakdash}{\nobreak\hbox{-}}
\setlength{\emergencystretch}{2em}
\multlinegap0pt
\usepackage{amssymb}
\usepackage{mathtools}
\newtheorem{lemma}{Lemma}
\title{Programs as Singularities}
\author{Daniel Murfet, Will Troiani}
\date{Source: arXiv:2504.08075v2 (CC BY 4.0)}
\begin{document}
\maketitle

\noindent\emph{Source and licence.} Programs as Singularities. Version arXiv:2504.08075v2 (CC BY 4.0), distributed by arXiv under the Creative Commons Attribution 4.0 International licence, http://creativecommons.org/licenses/by/4.0/, as recorded in the arXiv licence metadata for this exact version and shown on the abstract page https://arxiv.org/abs/2504.08075v2. Full licence notice: \texttt{licenses/arxiv-2504.08075-CC-BY-4.0.txt}.\par\smallskip

\noindent\textit{Editorial scope.} Counted unit: the article's lemma lemma:comparable\_hessian (source0.tex lines 2653--2663). The article's complete original proof of it is delivered with it (source0.tex lines 2665--2708, one \texttt{proof} environment of 44 lines). No mathematics is written, repaired or re-derived: the statement and the proof are the article's own lines, in the article's own notation, and no retained line is substituted or re-flowed.  No further source-local context is needed: every cross-reference of every retained line resolves inside the two delivered spans.  Nothing was omitted from any retained span, so the package carries no omission record.  No retained line contains a citation, so the wrapper carries no bibliography and no bibliography entry.  No retained line contains a figure environment, an image-inclusion command or drawing code, so no image is delivered and the package declares no asset dependency.  Editorial formatting only, in addition: an AMS \texttt{amsart} preamble that restates the article's own declarations and macros used by the retained lines, namely the environments \texttt{lemma} on separate counters, together with the standard packages those lines need.  The retained environments print in this document as Lemma 1 in order of appearance, whereas the article prints the same items with the numbering of its own full document.  The complete native source of the article as distributed in the arXiv e-print for this exact version is bundled unchanged as \texttt{sources/176342/source0.tex}.
\begin{lemma}\label{lemma:comparable_hessian}
Let $f, g: \mathbb{R}^d \to \mathbb{R}$ be real analytic functions on an open neighbourhood $U$ of the origin, with local minima at the origin and $f(0) = g(0) = 0$. Suppose there exist constants $c, d > 0$ such that 
\begin{equation}
c \cdot f(x) \leq g(x) \leq d \cdot f(x)
\end{equation}
for all $x \in U$. Let $Q_f$ and $Q_g$ be the Hessian matrices of $f$ and $g$ at the origin, respectively. Then:
\begin{enumerate}
\item[(i)] $\ker(Q_f) = \ker(Q_g)$.
\item[(ii)] For any unit eigenvector $v$ of $Q_f$ with eigenvalue $\lambda_f > 0$, we have $c \cdot \lambda_f \leq v^T Q_g v \leq d \cdot \lambda_f$.
\end{enumerate}
\end{lemma}

\begin{proof}
By hypothesis $f(0) = g(0) = 0$. Since $f$ and $g$ have local minima at $0$, we have $\nabla f(0) = \nabla g(0) = 0$, and both $Q_f$ and $Q_g$ are positive semidefinite. The Taylor expansions around the origin are:
\begin{align}
f(x) &= \frac{1}{2}x^T Q_f x + O(\|x\|^3)\\
g(x) &= \frac{1}{2}x^T Q_g x + O(\|x\|^3)
\end{align}
For (i) it suffices to show $\ker(Q_f) \subseteq \ker(Q_g)$. Let $v \in \ker(Q_f)$, so $Q_f v = 0$ and $v^T Q_f v = 0$. For sufficiently small $t$, consider $x = tv \in U$:
\begin{align}
f(tv) &= \frac{1}{2}t^2 v^T Q_f v + O(t^3) = O(t^3)
\end{align}

By our growth rate assumption, we have:
\begin{align}
g(tv) &\geq c \cdot f(tv) = c \cdot O(t^3) = O(t^3)\\
g(tv) &\leq d \cdot f(tv) = d \cdot O(t^3) = O(t^3)
\end{align}

However, we also know:
\begin{align}
g(tv) &= \frac{1}{2}t^2 v^T Q_g v + O(t^3)
\end{align}
For these equations to be consistent, we must have $v^T Q_g v = 0$. Since $Q_g$ is positive semidefinite, this implies $Q_g v = 0$, and therefore $v \in \ker(Q_g)$.

For (ii) let $v$ be a unit eigenvector of $Q_f$ with eigenvalue $\lambda_f > 0$, so $Q_f v = \lambda_f v$ and $v^T Q_f v = \lambda_f$.

For small $t$, consider $x = tv$:
\begin{align}
f(tv) &= \frac{1}{2}t^2 v^T Q_f v + O(t^3) = \frac{1}{2}t^2 \lambda_f + O(t^3)\\
g(tv) &= \frac{1}{2}t^2 v^T Q_g v + O(t^3)
\end{align}
From our growth rate constraint, we have:
\begin{align}
c \cdot f(tv) \leq g(tv) \leq d \cdot f(tv)
\end{align}
Substituting the expansions:
\begin{align}
c \cdot \left(\frac{1}{2}t^2 \lambda_f + O(t^3)\right) \leq \frac{1}{2}t^2 v^T Q_g v + O(t^3) \leq d \cdot \left(\frac{1}{2}t^2 \lambda_f + O(t^3)\right)
\end{align}
Dividing by $\frac{1}{2}t^2$ and taking the limit as $t \to 0$, we obtain:
\begin{align}
c \cdot \lambda_f \leq v^T Q_g v \leq d \cdot \lambda_f
\end{align}
This establishes the desired bounds on the quadratic form $v^T Q_g v$ for any unit eigenvector $v$ of $Q_f$ with positive eigenvalue $\lambda_f$.
\end{proof}

\end{document}
